% !TeX root = ../../Book3.tex
% 纲要标题：### 22.3 Gorenstein 环
% 纲要位置：计划纲要.md，第 2204 行。

\section{Gorenstein 环}
\label{b3:sec:2203}

环自身的有限内射消解，比所有模具有有限自由消解的要求宽松。一般 Bass 判别把这种有限性转为 Cohen–Macaulay 性与顶次 Ext 的一维条件，完全交则可由 Koszul 对偶满足这一条件。


% 纲要内容：
%
% * Noether 局部环的有限自内射维数
% * 与 Cohen–Macaulay 性的联系
% * 完全交蕴含 Gorenstein 的适当版本
%

\subsection{Noether 局部环的有限自内射维数}
\label{b3:subsec:2203:01}

正则局部环用有限自由消解控制所有模；Gorenstein 性则问环本身能否有有限的内射消解。两种要求的方向不同。特别是，剩余域有无限自由消解，并不排除环作为自身的模内射。

\begin{definition}\label{b3:def:injective-dimension-gorenstein}
设 \(R\) 为交换环，\(N\ne0\) 为 \(R\) 模。若存在正合列
\[
0\longrightarrow N\longrightarrow E^0\longrightarrow\cdots
\longrightarrow E^s\longrightarrow0
\]
且每个 \(E^i\) 内射，则满足此条件的最小 \(s\) 称为 \(N\) 的\term[neisheweishu]{内射维数}[injective dimension]，记为 \(\operatorname{id}_RN\)；若不存在有限长度的此种消解，则记为 \(\infty\)。对交换 Noether 局部环 \(R\)，条件 \(\operatorname{id}_RR<\infty\) 成立时，称 \(R\) 为 \term[Gorensteinhuan]{Gorenstein 环}[Gorenstein ring]。交换 Noether 环在每个素理想处的局部环均有此性质时，也称为 Gorenstein 环。
\end{definition}
\symidx{\(\operatorname{id}_RN\)：内射维数}

内射消解中的模通常不是有限模。譬如离散赋值环 \(R\) 的分式域 \(K\) 给出
\[
0\longrightarrow R\longrightarrow K\longrightarrow K/R\longrightarrow0.
\]
两个右端模都可除，由主理想整环上的 Baer 判据都内射；而 \(R\) 不是内射，因为均匀参数的倒数不在 \(R\) 中。因此 \(\operatorname{id}_RR=1\)。这里仍是用一个有限长度的消解描述一个一维环，却不能要求消解各项有限生成。

\begin{lemma}\label{b3:lem:injective-dimension-test}
设 \(R\) 为交换环，\(N\ne0\) 为 \(R\) 模，\(s\ge0\)。下列两条件等价：
\[
\operatorname{id}_RN\le s;\qquad
\operatorname{Ext}_R^{s+1}(R/J,N)=0
\text{ 对每个理想 }J\subseteq R\text{ 成立}.
\]
\end{lemma}
\begin{proof}
第一条件推出更高次的 Ext 全部为零。反过来，在一个内射消解中截取
\(0\to N\to E^0\to\cdots\to E^{s-1}\to C\to0\)；\(s=0\) 时令 \(C=N\)。逐次使用第二变量长正合列，得到
\(\operatorname{Ext}_R^1(R/J,C)\cong\operatorname{Ext}_R^{s+1}(R/J,N)=0\)。
由 \(0\to J\to R\to R/J\to0\) 的 Hom 长正合列，每个 \(J\to C\) 均可延拓到 \(R\)。第二卷第7.2节的 Baer 判据使 \(C\) 内射，故上述截断就是长度至多 \(s\) 的内射消解。

这里使用的内射计算与此前用自由消解定义的 Ext 一致：把自由消解 \(P_\bullet\to R/J\) 和内射消解 \(N\to E^\bullet\) 组成 \(\operatorname{Hom}_R(P_p,E^q)\)，两个方向分别由自由性和内射性保持正合，因而分别只剩 \(\operatorname{Hom}_R(P_\bullet,N)\) 与 \(\operatorname{Hom}_R(R/J,E^\bullet)\)。给定总次数中只有有限个项，双复形的圈消去论证给出相同上同调。
\end{proof}

\subsection{Gorenstein 性与 Cohen–Macaulay 性}
\label{b3:subsec:2203:02}

有限自内射维数不仅给出一项同调界，还强迫环达到 Cohen--Macaulay 深度。下面的判别不要求环完备，也不要求事先给定正则局部覆盖。证明采用素点间的 Ext 比较与正则元素换环，可对照松村英之\cite[\S 18, Lemmas 1--4 and Theorem 18.1, pp. 139--144]{Matsumura1989CommutativeRingTheory}。

\begin{theorem}[Gorenstein 环的判别]\label{b3:thm:gorenstein-bass-criterion}
设 \((R,\mathfrak m,k)\) 为交换 Noether 局部环，\(d=\dim R\)。下列条件等价：
\begin{equivalences}
\item \(R\) 为 Gorenstein 环；
\item \(R\) 为 Cohen--Macaulay 环，且
\(\dim_k\operatorname{Ext}_R^d(k,R)=1\)；
\item \(\operatorname{Ext}_R^i(k,R)=0\) 对 \(i\ne d\) 成立，且
\(\operatorname{Ext}_R^d(k,R)\cong k\)。
\end{equivalences}
这些条件成立时，\(\operatorname{id}_RR=d\)。
\end{theorem}
\begin{proof}
先说明三个计算事实。第一，\cref{b3:lem:injective-dimension-test}中的理想可只取素理想：每个有限模有有限素循环滤过，长正合列把素循环商上的 Ext 消失传到任意有限模，特别是各个 \(R/J\)。

第二，若 \(x\in\mathfrak m\) 是 \(R\) 非零因子，令 \(\overline R=R/xR\)。对每个 \(\overline R\) 模 \(N\)，有
\[
\operatorname{Hom}_R(N,R)=0,\qquad
\operatorname{Ext}_R^{j+1}(N,R)
\cong\operatorname{Ext}_{\overline R}^{j}(N,\overline R)
\quad(j\ge0).
\]
确实，二项自由消解 \(0\to R\xrightarrow{x}R\to\overline R\to0\) 给出
\(\operatorname{Ext}_R^q(\overline R,R)\) 只在 \(q=1\) 为 \(\overline R\)。把 \(N\) 的自由 \(\overline R\) 消解与 \(R\) 的内射消解组成 Hom 双复形，按\cref{b3:lem:canonical-change-of-rings}证明中的单行消去，便得上述换环同构；每个总次数只有有限项，不涉及无限乘积。

第三，对有限 \(R\) 模 \(M\)、素理想 \(\mathfrak p\) 及 \(t=\dim R/\mathfrak p\)，有
\[
\operatorname{Ext}_R^{j+t}(k,M)=0
\ \Longrightarrow\
\operatorname{Ext}_{R_{\mathfrak p}}^j
  (\kappa(\mathfrak p),M_{\mathfrak p})=0.
\]
先设 \(t=1\)，取 \(a\in\mathfrak m\setminus\mathfrak p\)。模 \(R/(\mathfrak p,a)\) 有有限长度，故给定的 \((j+1)\) 次消失也对它成立。由
\(0\to R/\mathfrak p\xrightarrow{a}R/\mathfrak p\to R/(\mathfrak p,a)\to0\)，乘以 \(a\) 在有限模 \(\operatorname{Ext}_R^j(R/\mathfrak p,M)\) 上满射；Nakayama 引理使该模为零，局部化即得断言。一般 \(t\) 时，取长度为 \(t\) 的最长素理想链，从 \(\mathfrak m\) 到 \(\mathfrak p\) 逐个作这一步；每个相邻商的局部维数为一，Ext 与有限模的局部化相容。

现在设 \(r=\operatorname{id}_RR<\infty\)。取满足 \(\dim R/\mathfrak p=d\) 的极小素理想，则零维局部环 \(R_{\mathfrak p}\) 有非零基座。第三个事实的逆否命题给出 \(\operatorname{Ext}_R^d(k,R)\ne0\)，所以 \(r\ge d\)。若 \(r>0\)，则 \(\operatorname{Ext}_R^r(k,R)\ne0\)：第一事实给出某个 \(\mathfrak p\) 使 \(\operatorname{Ext}_R^r(R/\mathfrak p,R)\ne0\)；若 \(\mathfrak p\ne\mathfrak m\)，取 \(a\in\mathfrak m\setminus\mathfrak p\)，同一短正合列及 \((r+1)\) 次 Ext 消失会使乘以 \(a\) 满射，再由 Nakayama 引理得到矛盾。

此时 \(\mathfrak m\notin\operatorname{Ass}R\)。否则有单射 \(k\to R\)，其余核的 \((r+1)\) 次 Ext 为零，迫使
\(0=\operatorname{Ext}_R^r(R,R)\to\operatorname{Ext}_R^r(k,R)\) 满射，矛盾。因此可选 \(x\in\mathfrak m\) 为非零因子。第二事实和第一事实使 \(\operatorname{id}_{\overline R}\overline R=r-1\)：上界由全部素循环商的换环同构得到，下界由 \(\operatorname{Ext}_R^r(k,R)\ne0\) 得到。对 \(r\) 归纳，\(r=0\) 时已有 \(d=0\)，得到 \(R\) 为 CM 环且 \(r=d\)。

归纳还给出条件3。零维且自内射时，任取单射 \(k\to R\)，内射性使 \(\operatorname{Hom}_R(R,R)\to\operatorname{Hom}_R(k,R)\) 满射。后者由该单射生成且被 \(\mathfrak m\) 零化，故同构于 \(k\)，正次 Ext 全为零。正维时，\(\operatorname{Hom}_R(k,R)=0\)，其余各次由第二事实下降到 \(\overline R\)。这证明条件1推出条件3；条件3结合深度的 Ext 刻画显然推出条件2。

最后设条件2成立。取参数正则序列 \(x_1,\ldots,x_d\)，令 \(T=R/(x_1,\ldots,x_d)\)。反复换环得到
\(\operatorname{Soc}(T)\cong\operatorname{Ext}_R^d(k,R)\cong k\)。Artin 局部环中每个非零理想都与基座非零相交：对该理想反复乘以极大理想，取最后一个非零幂次即可。将基座同构 \(\operatorname{Soc}(T)\to k\subseteq E_T(k)\) 延拓到 \(T\to E_T(k)\)，本质性使其为单射；由\cref{b3:lem:matlis-finite-length}，
\(\ell_T E_T(k)=\ell_T D_T(T)=\ell_T T\)，故它是同构，\(T\) 自内射。

沿正则序列逐项返回。若 \(\operatorname{id}_{\overline R}\overline R\le s\)，则对含 \(x\) 的素理想，换环给出 \((s+2)\) 次 Ext 消失；对不含 \(x\) 的素理想，模 \(R/(\mathfrak p,x)\) 被 \(x\) 零化，短正合列使乘以 \(x\) 在 \(\operatorname{Ext}_R^{s+2}(R/\mathfrak p,R)\) 上满射，Nakayama 引理再次使其为零。第一事实于是给出 \(\operatorname{id}_RR\le s+1\)。从 \(T\) 返回 \(R\) 得 \(\operatorname{id}_RR\le d\)，而顶次 Ext 非零给出等号，完成证明。
\end{proof}

对一个 \(d\) 维 CM 局部环 \(R\)，整数
\(\dim_k\operatorname{Ext}_R^d(k,R)\) 称为它的\term[CohenMacaulayleixing]{Cohen–Macaulay 类型}[Cohen--Macaulay type]。深度的 Ext 刻画保证这个整数为正。判别定理说明 Gorenstein 环恰是类型为一的 CM 环；只有“CM”仍不足够。零维环的类型就是其基座的维数，所以 \(k[x,y]/(x,y)^2\) 是 CM 环但不是 Gorenstein 环。

\begin{proposition}\label{b3:prop:gorenstein-completion-localization}
设 \((R,\mathfrak m,k)\) 为交换 Noether 局部环。则 \(R\) 为 Gorenstein 环当且仅当其 \(\mathfrak m\)-进完备化 \(\widehat R\) 为 Gorenstein 环。若 \(R\) 为 Gorenstein 环，则每个 \(R_{\mathfrak p}\)、\(\mathfrak p\in\operatorname{Spec}R\)，也为 Gorenstein 环。
\end{proposition}
\begin{proof}
完备化保维数并忠实平坦，且由\cref{b3:thm:cm-completion}保持及反映 CM 性。取 \(k\) 的有限秩自由 \(R\) 消解，逐项张量 \(\widehat R\) 后仍是 \(k\) 的自由消解。有限秩保证 Hom 与此张量相容，平坦性保证上同调与张量相容，故
\[
\operatorname{Ext}_R^i(k,R)\otimes_R\widehat R
\cong\operatorname{Ext}_{\widehat R}^i(k,\widehat R).
\]
左侧原本被 \(\mathfrak m\) 零化，张量后不变。应用上一判别定理，便得完备化等价性。

为处理局部化，先说明内射模 \(E\) 的局部化 \(E_{\mathfrak p}\) 对 \(R_{\mathfrak p}\) 内射。该局部环的每个理想都是某个有限生成理想 \(J\subseteq R\) 的局部化。Noether 性使 \(J\) 有限展示，从而
\[
\operatorname{Hom}_{R_{\mathfrak p}}(J_{\mathfrak p},E_{\mathfrak p})
\cong\operatorname{Hom}_R(J,E)_{\mathfrak p}.
\]
原内射性使 \(\operatorname{Hom}_R(R,E)\to\operatorname{Hom}_R(J,E)\) 满射，局部化后仍满射，再用 Baer 判据即可。把 \(R\) 的有限内射消解逐项局部化，就得到 \(R_{\mathfrak p}\) 的有限内射消解。
\end{proof}

在给定正则覆盖的情形，典范模提供另一种直接的充分条件，并说明自内射维数为何能由局部对偶控制。

\begin{proposition}\label{b3:prop:canonical-free-gorenstein}
设 \((S,\mathfrak n)\) 为 \(n\) 维交换 Noether 完备正则局部环，\(I\subsetneq S\) 为理想，\(R=S/I\) 为 \(d\) 维 Cohen--Macaulay 环，典范模为
\(\omega_R=\operatorname{Ext}_S^{n-d}(R,S)\)。若 \(\omega_R\cong R\) 为 \(R\) 模，则 \(R\) 为 Gorenstein 环。
\end{proposition}
\begin{proof}
对每个理想 \(J\subseteq R\)，模 \(R/J\) 有限。\cref{b3:thm:canonical-module-local-duality}给出
\[
\operatorname{Ext}_R^{d+1}(R/J,\omega_R)
\cong D_R\bigl(H_{\mathfrak m}^{-1}(R/J)\bigr)=0.
\]
用给定的 \(\omega_R\cong R\) 替换第二变量，再由\cref{b3:lem:injective-dimension-test}得 \(\operatorname{id}_RR\le d\)。此外 CM 性及\cref{b3:thm:depth-ext}使 \(\operatorname{Ext}_R^d(k,R)\ne0\)，所以内射维数恰为 \(d\)。
\end{proof}

此处典范模同构是一个 \(R\) 模同构，通常没有指定的自然选择。把典范模选成 \(R\) 也不等于对每个有限模使用普通线性对偶 \(\operatorname{Hom}_R(-,R)\)：局部对偶中的次数 \(d-i\) 和 Matlis 对偶仍然不可省略。

\subsection{完全交的 Gorenstein 性}
\label{b3:subsec:2203:03}

完全交商的 Koszul 消解有一种对称性：取 \(S\) 线性对偶后，复形各项的顺序倒过来，微分仍是原正则序列给出的矩阵。这使典范模的计算十分具体。

\begin{proposition}\label{b3:prop:ci-canonical-module}
设 \((S,\mathfrak n)\) 为 \(n\) 维交换 Noether 完备正则局部环，\(f_1,\ldots,f_c\in\mathfrak n\) 为 \(S\) 正则序列，\(R=S/(f_1,\ldots,f_c)\)。则 \(d=\dim R=n-c\)，而且
\[
\operatorname{Ext}_S^j(R,S)=
\begin{cases}
R,&j=c,\\
0,&j\ne c.
\end{cases}
\]
特别地，\(R\) 的典范模 \(\omega_R\) 同构于 \(R\)。
\end{proposition}
\begin{proof}
\cref{b3:thm:complete-intersection-cm}及\cref{b3:prop:given-complete-intersection}给出 CM 性和维数。用\cref{b3:prop:ci-koszul-resolution}的 Koszul 消解 \(K_\bullet(f;S)\) 计算 Ext。外积给出完美配对
\[
\bigwedge^jS^c\times\bigwedge^{c-j}S^c
\longrightarrow\bigwedge^cS^c\cong S,
\]
最后一个同构取有序基 \(e_1\wedge\cdots\wedge e_c\mapsto1\)。若 \(u\) 的次数为 \(j\)，Koszul 微分满足
\(\mathrm{d}(u\wedge v)=\mathrm{d}u\wedge v+(-1)^ju\wedge \mathrm{d}v\)。因此配对使对偶微分与倒序 Koszul 微分相差一个只依赖次数的符号；逐次把某些次数上的配对乘以 \(-1\)，便得到复形同构
\[
\operatorname{Hom}_S(K_\bullet(f;S),S)^j
\cong K_{c-j}(f;S).
\]
正则序列使右侧只有链次数零的同调 \(R\)，所以对偶上同调只在 \(j=c\) 为 \(R\)。\(c=0\) 时消解只有 \(S\)，结论同样成立。由于 \(n-d=c\)，这正是典范模定义要求的次数。
\end{proof}

\begin{theorem}\label{b3:thm:complete-intersection-gorenstein}
每个交换 Noether 完全交局部环都是 Gorenstein 环。因此，对交换 Noether 局部环有
\[
\text{正则}\ \Longrightarrow\ \text{完全交}\
\Longrightarrow\ \text{Gorenstein}\
\Longrightarrow\ \text{Cohen--Macaulay}.
\]
\end{theorem}
\begin{proof}
按照\cref{b3:def:local-complete-intersection}，完全交局部环 \(R\) 的完备化具有
\(\widehat R=S/(f_1,\ldots,f_c)\) 的表示，其中 \(S\) 完备正则局部，\(f_1,\ldots,f_c\) 为正则序列。前一命题使 \(\omega_{\widehat R}\cong\widehat R\)，再由\cref{b3:prop:canonical-free-gorenstein}得 \(\widehat R\) 为 Gorenstein 环。\cref{b3:prop:gorenstein-completion-localization}把此性质下降到 \(R\)。链中的第一项由\cref{b3:thm:complete-intersection-cm}证明中的空正则序列论证成立，最后一项来自\cref{b3:thm:gorenstein-bass-criterion}。
\end{proof}

这里没有为任意完备局部环暗中选择一个正则局部覆盖：完全交的定义已经提供了本次所需覆盖。一般完备局部环的覆盖存在性属于 Cohen 结构定理的范围，安排在本卷附录。对于一个已经给定的正则局部环 \(S\) 及非零非单位 \(f\)，商 \(S/(f)\) 是超曲面，因而也是 Gorenstein 环；它是否正则仍须另外判断。
